\section{Results}

\subsection{Exact null distributions replace permutation testing}
For every gene set, NetFunrankR computes the null distributions of the overlap, linker count and distance statistics exactly (Figure~\ref{fig:schematic}). Exact p-values agreed with 2{,}000 degree-binned Monte-Carlo permutations (Spearman $\rho = 0.9994$ between $\log_{10}$ p-values of the 320 KEGG pathways for the COPD list; all differences within Monte-Carlo error), and the saddlepoint approximation used for long lists differed from the exact values by at most $0.02$ $\log_{10}$ units (Supplementary Figure~S1). Unlike permutation p-values, which cannot fall below $1/(B+1)$, the exact p-values resolve the extreme tail (e.g.\ $p = 1.3\times10^{-7}$ for \textit{Nicotine addiction} in the COPD analysis, where 2{,}000 permutations give $p \ge 5\times10^{-4}$). Removing the permutation step also removes the main computational bottleneck (Supplementary Figure~S2): the original NetFunrankR v1 implementation required 26\,s per pathway for only 100 permutations (2.3\,h for the 320 KEGG pathways), whereas v2 tests all KEGG pathways in 0.5--3.4\,s and all Reactome pathways in 3.4--24\,s for query lists of 10--2{,}000 genes on a single laptop core, after a one-off distance computation per library (2.8\,s for KEGG, 20\,s for Reactome). This is comparable to Proximity-$z$ with 1{,}000 permutations (0.1--20\,s), whose p-values, however, rest on a normal approximation of the permutation $z$-score that is miscalibrated (Section~\ref{res:calibration}).

\subsection{Hub-biased gene lists inflate false discoveries of ORA and published network methods}\label{res:calibration}
We first asked whether each method controls false positives when the query list contains no biological signal (Figure~\ref{fig:calibration}). For uniformly sampled random lists, ORA was conservative (1.6--3.3\% of pathways at $p<0.05$), whereas NEA was strongly anti-conservative (15.6--19.2\%), and 100\% of random lists produced at least one FDR-significant pathway with NEA. Random lists that over-represent highly connected proteins---as disease-gene and GWAS lists do---changed the picture for ORA: its per-pathway false-positive rate rose with list size to 7.5\% ($n = 200$), and the fraction of random lists with at least one FDR-significant pathway reached 26\%, 74\%, 99\% and 100\% for $n = 20$, 50, 100 and 200, respectively. Permutation-based proximity with a normal approximation (Proximity-$z$; NetFunrankR v1 statistic) was mildly anti-conservative (per-pathway rates up to 6.3\% and 6.9\%), yielding spurious FDR-significant pathways for 5--13\% and 9--22\% of random lists. All NetFunrankR v2 tests were calibrated in every setting: per-pathway false-positive rates were at most 5.4\% (combined test: 3.2--4.9\%), and at most 5.5\% of random lists produced any FDR-significant pathway with the combined test. RWR with 200 permutations was calibrated but, owing to its permutation floor, could not produce p-values below $5\times10^{-3}$ (flat segment in Figure~\ref{fig:calibration}a), which made FDR-significant findings practically unreachable (Sections~\ref{res:power}--\ref{res:gwas}). These results show that the degree-preserving null is not a technicality: without it, hub-rich gene lists generate pathway ``enrichments'' by construction.

\subsection{Network proximity recovers pathways with incomplete annotation}\label{res:power}
We next planted pathway signal into random query lists while hiding a growing fraction of the pathway's genes from the tested annotation (Figure~\ref{fig:power}; 100 KEGG pathways, 50-gene lists). With complete annotation all overlap-based tests detected the target pathway when 20\% of the query came from it (power 0.99--1.00), whereas distance-only proximity (Proximity-$z$ 0.81; NetFunrankR total proximity 0.79) and RWR (0.41) were less sensitive. As annotation became incomplete, ORA lost power rapidly (0.54 and 0.24 with 75\% and 90\% of the pathway hidden), while the NetFunrankR combined test retained 0.80 and 0.62, because hidden pathway members were recovered as linker genes. The combined test also ranked the target pathway best among calibrated methods (median rank percentile 1.3\% and 3.5\% at 75\% and 90\% hidden, versus 6.1\% and 15.4\% for ORA). NEA (0.93 and 0.73) and the NetFunrankR v1 statistic (0.77 and 0.49) were also robust, but at the price of the type-I error inflation shown above; NEA ranked the target lower than NetFunrankR in every setting. With strong signal (40\%) every method except RWR exceeded power 0.85 at 90\% hidden annotation, and NetFunrankR reached 0.93 versus 0.58 for ORA. The cost of combining two tests appeared only for the weakest signal (10\%, i.e.\ five pathway genes) with complete annotation, where the combined test had lower power than ORA (0.70 vs 0.90) because the linker component carried no information; with 25\% or more of the annotation missing the two were equivalent.

\subsection{Expression benchmark (GSEABenchmarkeR)}\label{res:gseabench}
On the 42 GEO2KEGG expression datasets (Figure~\ref{fig:gseabench}), gene-list methods applied to the differentially expressed genes ranked MalaCards-relevant pathways at least as well as methods using the full expression data, in line with the original benchmark, in which ORA and PADOG were among the best performers \citep{Geistlinger2021}. NetFunrankR (64.8\% of the optimal relevance score) outperformed ORA (60.4\%; better in 27/42 datasets, paired Wilcoxon $p = 0.006$), PADOG (61.5\%; 29/42, $p = 0.045$), SAFE (58.1\%; $p = 6\times10^{-4}$), GSEA (54.5\%; 35/42, $p = 1.6\times10^{-6}$), CAMERA (53.8\%; $p = 7\times10^{-5}$), GSVA, ROAST, SAMGS and RWR (61.2\%; 35/42, $p = 8\times10^{-5}$), and did not differ significantly from GlobalTest (61.8\%; $p = 0.08$), NEA (64.9\%; $p = 0.79$) or Proximity-$z$ (66.9\%; $p = 0.07$). The degree-aware overlap test alone again improved on the hypergeometric test (28/42 datasets, $p = 0.007$ for the combined versus overlap-only comparison; Figure~\ref{fig:gseabench}a). The number of significant sets distinguished the network methods sharply: NEA and Proximity-$z$ declared a median of 44\% and 74\% of all KEGG pathways significant per dataset---as did GlobalTest (89\%) and SAMGS---whereas NetFunrankR reported a median of 8.1\%, comparable to CAMERA (8.7\%) and GSVA (7.0\%). With long expression-derived lists (median 1{,}845 DE genes), NetFunrankR used the saddlepoint approximation and required a median of 11\,s per dataset, compared with 70\,s for PADOG and 176\,s for GSEA with 1{,}000 sample permutations.

\subsection{GWAS gene lists: three independent gold standards}\label{res:gwas}
Gene lists from GWAS are the setting in which annotation-based enrichment is weakest, because many associated genes are not (yet) annotated to the pathways through which they act. We assembled genome-wide significant gene lists for 25 diseases (15--1{,}387 network genes; Supplementary Table~S1) and evaluated all methods against three gold standards (Figure~\ref{fig:gwas}).

\emph{MalaCards relevance.} All network methods ranked phenotype-relevant pathways higher than ORA (mean relevance 59.1\% of the optimum). The NetFunrankR combined test reached 74.4\%, outperforming ORA in 21 of 25 diseases (paired Wilcoxon $p = 3.2\times10^{-5}$), RWR (69.9\%; 21/25, $p = 2.2\times10^{-5}$), Proximity-$z$ (69.6\%; 20/25, $p = 6.3\times10^{-4}$) and the distance-based combination (71.2\%; 17/25, $p = 0.008$), and performing on par with NEA (75.0\%; $p = 0.56$) and the NetFunrankR v1 statistic (73.0\%; $p = 0.21$)---the two methods with inflated type-I error. Notably, the degree-aware overlap test alone improved on the hypergeometric test (66.3\% vs 59.1\%; 21/25 diseases, $p = 4.2\times10^{-6}$), showing that correcting hub bias also improves the \emph{ranking} of relevant pathways.

\emph{Significant findings.} Because ranking metrics ignore calibration, we also examined what users act on: pathways called at FDR $<0.05$. The combined test reported 1{,}259 significant disease--pathway pairs across 25 diseases, of which 321 were MalaCards-relevant, compared with 148 relevant among 438 for ORA; i.e.\ NetFunrankR more than doubled the number of relevant significant pathways and produced significant results for 22 rather than 19 diseases, at a lower precision (25\% vs 34\%). Since the MalaCards rankings are themselves derived from gene overlap between disease genes and pathways, this precision estimate is expected to favour overlap-based methods. NEA reported 3{,}734 significant pairs, which, given its calibration (Section~\ref{res:calibration}), cannot be interpreted as discoveries.

\emph{Target disease pathway.} For the 21 diseases with a dedicated KEGG disease pathway, the combined test called the target pathway significant (FDR $<0.05$) in 11 diseases, compared with 8 for ORA and 3 for RWR; median target ranks were similar for ORA (7th percentile), the degree-aware overlap test (6th) and the combined test (9th).

\emph{Temporal validation.} When methods were applied only to the genes reported up to each disease's median discovery date, the pathways they prioritised predicted the pathways over-represented among genes discovered later (mean AUROC: NEA 0.795, NetFunrankR 0.767, RWR 0.757, ORA 0.755; 10 diseases). The difference between NetFunrankR and ORA was not significant (6/10 diseases; $p = 0.70$), whereas both distance-only statistics (Proximity-$z$ 0.704; NetFunrankR distance test 0.641) performed worst, confirming that the linker formulation is the better beyond-overlap statistic for sparse genetic signals.

\subsection{Robustness}\label{res:robust}
The GWAS benchmark was repeated under alternative interactomes and perturbations (Figure~\ref{fig:robust}). Removing STRING's curated-database and text-mining channels did not cost performance---it improved it: the combined test reached 74.4\% relevance on the non-circular network versus 71.9\% on the conventional full-evidence network (score $\ge 0.7$) and 71.7\% on the smaller high-confidence non-circular network. Results were insensitive to the degree-bin size (74.2\% for minimum bin sizes of 50 and 250) and to missing interactions (73.9\% and 73.7\% after removing 25\% and 50\% of the edges). As a negative control, degree-preserving rewiring of the interactome, which keeps every protein's degree but destroys biological wiring, reduced the linker test from 70.5\% to 58.4\%---the level of ORA---and the distance test from 64.7\% to 53.2\%, while the overlap component was unaffected. Hence the beyond-overlap evidence reflects the specific topology of the interactome rather than degree artefacts. Finally, results were as stable as ORA under subsampling of the query: the top-20 pathways obtained from random halves of each GWAS list overlapped those from the full list with a mean Jaccard index of 0.44 for the combined test and 0.47 for ORA, compared with 0.40 for Proximity-$z$ and 0.36 for RWR.

\subsection{Case study: COPD}\label{res:copd}
The 307 COPD-associated genes of \citet{MorselliGysi2021} (261 in the network) were tested against 363 KEGG pathways (Figure~\ref{fig:copd}). The combined test identified 24 pathways at FDR $<0.05$: 11 supported mainly by overlap, 7 by both overlap and linker genes, and 6 by network evidence only, compared with 7 pathways for the hypergeometric ORA. The NetFunrankR ranking agreed best with MalaCards COPD relevance (77.9\% of optimum; NEA 73.9\%, ORA 70.0\%, Proximity-$z$ 68.1\%, RWR 66.4\%). It was also the ranking most consistent with an independent \emph{expression} analysis of COPD (CAMERA on peripheral blood transcriptomes, GSE42057; Spearman $\rho = 0.10$ for the combined and $0.14$ for the linker test, versus $0.00$ for ORA), although agreement between genetic and transcriptomic evidence was weak for all methods. The advantage of network evidence was largest for short lists: with random subsets of 6 or 10 COPD genes, ORA recovered only 9.4\% and 8.5\% of the optimal relevance score, whereas NetFunrankR reached 54.9\% and 52.8\%, comparable to the other network methods but with calibrated significance (Figure~\ref{fig:copd}c).
